\documentclass[12pt]{article}
\usepackage{amsmath, amssymb, amsthm}
\usepackage{geometry}
\usepackage{hyperref}
\usepackage{bm}

\geometry{margin=1in}

\newcommand{\be}{\begin{equation}}
\newcommand{\ee}{\end{equation}}
\newcommand{\bea}{\begin{aligned}}
\newcommand{\eea}{\end{aligned}}

\theoremstyle{plain}
\newtheorem{claim}{Claim}

\title{\textbf{Coupled Harmonic Oscillators on a Ring}}
\date{}

\begin{document}
\maketitle

\noindent\textbf{Problem 2.} (Coupled harmonic oscillators)
Consider a linear system of $N$ particles coupled by springs, governed by the Hamiltonian
\be
  H = \frac{1}{2}\sum_{i=1}^N p_i^2 + \frac{\omega^2}{2}\sum_{i=1}^N (q_i - q_{i+1})^2,
  \qquad q_{N+1} \coloneqq q_1.
\ee
The Poisson brackets for the canonical variables are
\be
  \{q_k, q_l\} = 0, \qquad \{p_k, p_l\} = 0, \qquad \{p_k, q_l\} = -\delta_{k,l}.
\ee

\bigskip

%------------------------------------------------------------
\noindent\textbf{Part 1.} \textit{Hamiltonian equations, matrix $A$, its eigenvalues and eigenvectors.}

\medskip
\noindent The Hamiltonian equations read
\be
  \dot{q}_k = \{q_k, H\} = \frac{\partial H}{\partial p_k} = p_k, \qquad
  \dot{p}_k = \{p_k, H\} = -\frac{\partial H}{\partial q_k}
            = -\omega^2\bigl(2q_i - q_{i-1} - q_{i+1}\bigr).
\ee
In vector form these become
\be
  \dot{\bm{q}} = \bm{p}, \qquad \dot{\bm{p}} = -\omega^2 A\,\bm{q},
\ee
where the matrix $A$ encodes the nearest-neighbour coupling with periodic boundary conditions:
\be
  A = \begin{pmatrix}
        2  & -1 &  0 & \cdots & \cdots & -1 \\
       -1  &  2 & -1 &  0    & \cdots &  0 \\
        0  & -1 &  2 & -1    & \cdots &  0 \\
       \vdots & \vdots & & \ddots & & \vdots \\
       \vdots & \vdots & & & \ddots & \vdots \\
       -1  &  0 & \cdots & \cdots & -1 &  2
      \end{pmatrix}.
\ee

\noindent\textbf{Eigenvalues and eigenvectors.}
The matrix $A$ commutes with the cyclic permutation (translation) operator
\be
  T = \begin{pmatrix}
       0 & 1 & 0 & \cdots & 0 \\
       0 & 0 & 1 & \cdots & 0 \\
       \vdots & & \ddots & & \vdots \\
       \vdots & & & \ddots & \vdots \\
       1 & 0 & 0 & \cdots & 0
      \end{pmatrix},
\ee
so $A$ and $T$ share a common eigenbasis. Since $T^N = \mathbf{1}$, the eigenvalues of $T$ satisfy $\lambda_k^N = 1$, giving
\be
  \lambda_k = e^{i2\pi k/N}, \qquad k = 0, 1, \ldots, N-1.
\ee
The corresponding eigenvectors of $T$ (and hence of $A$) are the discrete Fourier modes
\be
  \bm{x}_k = \bigl(1,\; e^{i2\pi k/N},\; \ldots,\; e^{i2\pi k(N-1)/N}\bigr)^{\!\top}.
\ee
Applying $A$ to $\bm{x}_k$, the $j$-th component gives
\be\bea
  (A\bm{x}_k)_j
    &= e^{i2\pi kj/N}\bigl(2 - e^{i2\pi k/N} - e^{-i2\pi k/N}\bigr) \\
    &= 2e^{i2\pi kj/N}(1 - \cos(2\pi k/N)) \\
    &= 4\sin^2\!\Bigl(\frac{\pi k}{N}\Bigr)\,(\bm{x}_k)_j.
\eea\ee
Therefore the eigenvalues of $A$ are
\be
  \mu_k = 4\sin^2\!\Bigl(\frac{\pi k}{N}\Bigr), \qquad k = 0, 1, \ldots, N-1.
\ee
Note that $\mu_0 = 0$ (the zero mode corresponding to uniform translation) and $\mu_k = \mu_{N-k}$
(doubly degenerate for $k \neq 0, N/2$).

\bigskip

%------------------------------------------------------------
\noindent\textbf{Part 2.} \textit{Normal coordinates, Hamiltonian, and solution.}

\medskip
\noindent Set $N = 2M+1$ and define the discrete Fourier transforms
\be
  Q_k = \frac{1}{\sqrt{N}}\sum_{i=1}^N e^{\frac{2\pi i}{N}ki}\,q_i, \qquad
  P_k = \frac{1}{\sqrt{N}}\sum_{i=1}^N e^{\frac{2\pi i}{N}ki}\,p_i,
  \qquad k = -M, \ldots, M.
\ee
The inverse transformation is
\be
  q_j = \frac{1}{\sqrt{N}}\sum_{k=-M}^M e^{\frac{-2\pi i}{N}kj}\,Q_k, \qquad
  p_j = \frac{1}{\sqrt{N}}\sum_{k=-M}^M e^{\frac{-2\pi i}{N}kj}\,P_k.
\ee

\noindent\textbf{Hamiltonian in normal coordinates.}
Using the discrete orthogonality relation $\frac{1}{N}\sum_j e^{i2\pi(k+m)j/N} = \delta_{k,-m}$, we compute
\be\bea
  \sum_j p_j^2
    &= \frac{1}{N}\sum_{k,m,j} P_k P_m\,e^{\frac{i2\pi(k+m)j}{N}}
     = \sum_k P_{-k}P_k = \sum_k |P_k|^2, \\[6pt]
  \sum_{j=1}^N(q_j - q_{j+1})^2
    &= \frac{1}{N}\sum_{k,m,j}Q_k Q_m\,
       e^{-i\pi(k+m)(2j+1)/N}\,
       \bigl(2i\sin\tfrac{\pi k}{N}\bigr)
       \bigl(2i\sin\tfrac{\pi m}{N}\bigr) \\
    &= 4\sum_k |Q_k|^2 \sin^2\!\Bigl(\frac{\pi k}{N}\Bigr).
\eea\ee
Hence the Hamiltonian becomes a sum of decoupled oscillators:
\be
  \boxed{H = \frac{1}{2}\sum_{k=-M}^M |P_k|^2
             + 2\omega^2\sum_{k=-M}^M |Q_k|^2\sin^2\!\Bigl(\frac{\pi k}{N}\Bigr).}
\ee

\noindent\textbf{Equations of motion and their solution.}
The Hamiltonian equations in normal coordinates are
\be
  \dot{Q}_k = P_k, \qquad
  \dot{P}_k = -4\omega^2\sin^2\!\Bigl(\frac{\pi k}{N}\Bigr)Q_k,
\ee
which combine into $\ddot{Q}_k = -\omega(k)^2 Q_k$ with dispersion relation
\be
  \omega(k) = 2\omega\sin\frac{\pi k}{N}.
\ee
The general solution is
\be
  Q_k(t) = Ae^{i\omega(k)t} + Be^{-i\omega(k)t}, \qquad
  P_k(t) = i\omega(k)\bigl(Ae^{i\omega(k)t} - Be^{-i\omega(k)t}\bigr).
\ee
The zero mode ($k=0$) satisfies $\dot{P}_0 = 0$, so $Q_0(t) = Q_0(0) + P_0 t$ describes free
centre-of-mass motion.

\bigskip

%------------------------------------------------------------
\noindent\textbf{Part 3.} \textit{Creation/annihilation variables, conserved quantities, involutivity.}

\medskip
\noindent For $k = \pm 1, \ldots, \pm M$, define
\be
  a_k = \frac{1}{\sqrt{2}}\bigl(P_{-k} + i\omega(k)Q_{-k}\bigr), \qquad
  \bar{a}_k = \frac{1}{\sqrt{2}}\bigl(P_k + i\omega(k)Q_k\bigr).
\ee

\noindent\textbf{Hamiltonian in terms of $a_k$, $\bar{a}_k$.}
A direct computation gives
\be\bea
  a_{-k}a_k &= \tfrac{1}{2}\bigl(|P_k|^2 + \omega(k)^2|Q_k|^2
                + i\omega(k)(P_kQ_{-k} - P_{-k}Q_k)\bigr), \\
  \bar{a}_{-k}\bar{a}_k &= \tfrac{1}{2}\bigl(|P_k|^2 + \omega(k)^2|Q_k|^2
                + i\omega(k)(P_{-k}Q_k - P_kQ_{-k})\bigr),
\eea\ee
so that $a_{-k}a_k + \bar{a}_{-k}\bar{a}_k = |P_k|^2 + \omega(k)^2|Q_k|^2$. 
Since $2\omega^2\sin^2(\pi k/N) = \frac{1}{2}\omega(k)^2$, the Hamiltonian becomes
\be
  \boxed{H = \frac{1}{2}P_0^2 + \sum_{k=1}^M (a_{-k}a_k + \bar{a}_{-k}\bar{a}_k)
           = I_0 + \sum_{k=1}^M (I_k + I_{-k}).}
\ee

\noindent\textbf{Poisson brackets.}
From $\{P_k, Q_l\} = -\frac{1}{N}\sum_j e^{i2\pi j(k+l)/N} = -\delta_{k,-l}/N$
we derive the fundamental brackets of the normal coordinates, and from those:
\be\bea
  \{a_k, a_l\} &= \frac{i\omega(k)}{N}\,\delta_{k,-l}, \\[4pt]
  \{\bar{a}_k, \bar{a}_l\} &= \frac{i\omega(k)}{N}\,\delta_{k,-l}, \\[4pt]
  \{a_k, \bar{a}_l\} &= 0.
\eea\ee
The last bracket vanishes because the cross terms cancel:
\be
  2\{a_k, \bar{a}_l\}
  = i\omega(k)\{Q_{-k}, P_l\} + i\omega(l)\{P_{-k}, Q_l\}
  = \frac{i\omega(k)}{N}\delta_{k,l} - \frac{i\omega(l)}{N}\delta_{k,l} = 0.
\ee

\noindent\textbf{Conservation laws.}
We show $\dot{I}_0 = \dot{I}_k = \dot{I}_{-k} = 0$ by verifying that $I_0$, $I_k$, $I_{-k}$
each Poisson-commute with $H$.

\begin{claim}
$\{I_k, I_l\} = 0$ for all $k \neq l$, where $I_k = a_{-k}a_k$ and $I_{-k} = \bar{a}_{-k}\bar{a}_k$.
\end{claim}
\begin{proof}
Using the Leibniz rule and the bracket relations computed above:
\be\bea
  \{a_{-k}a_k,\; a_{-l}a_l\}
    &= a_{-k}\{a_k, a_{-l}\}a_l + a_{-k}a_{-l}\{a_k, a_l\}
     + \{a_{-k}, a_{-l}\}a_ka_l + a_{-l}\{a_{-k}, a_l\}a_k \\
    &= \frac{i\omega(k)a_{-k}a_k}{N}(\delta_{k,-l} + \delta_{k,l} - \delta_{k,l} - \delta_{k,-l})
     = 0, \\[4pt]
  \{\bar{a}_{-k}\bar{a}_k,\; \bar{a}_{-l}\bar{a}_l\}
    &= 0 \quad \text{(same argument)}, \\[4pt]
  \{a_{-k}a_k,\; \bar{a}_{-l}\bar{a}_l\}
    &= 0 \quad \text{(since $\{a_\bullet, \bar{a}_\bullet\} = 0$)}.
\eea\ee
\end{proof}

\noindent Since $H = I_0 + \sum_{k=1}^M(I_k + I_{-k})$ and $I_0 = \frac{1}{2}P_0^2$ trivially commutes
with all $I_k$, we conclude
\be
  \{I_0, H\} = 0, \qquad
  \{I_k, H\} = \sum_{l=1}^M\{I_k, I_l + I_{-l}\} = 0, \qquad
  \{I_{-k}, H\} = 0,
\ee
hence $\dot{I}_0 = \dot{I}_k = \dot{I}_{-k} = 0$ for all $k = 1, \ldots, M$.

\medskip
\noindent The system therefore possesses $2M + 1 = N$ independent conserved quantities
$\{I_0, I_1, \ldots, I_M, I_{-1}, \ldots, I_{-M}\}$ which are pairwise in involution,
confirming complete integrability of the coupled oscillator chain.

\end{document}
